\section{Discussion}
\label{sec:discussion}

\subsection{Why hierarchical pruning works here}
The ablation shows that structural restrictions, feasibility bounds and an incumbent
Pareto archive provide complementary reductions on these design spaces. The default
search does not need early queueing contributions to obtain the reported evaluation
savings. The separate order experiment (Section~\ref{sec:order}) measures the effect of
assigning their determinants earlier; it also makes clear that changing order changes
which incumbents become available for dominance pruning.

\subsection{When HCA-DSE is the right tool --- and when it is not}
The method is advantageous when the design space is large, the constraints are tight, the
evaluator is expensive, and admissible bounds can be constructed. The crossover analysis quantifies the third condition using actual exhaustive evaluator
calls and measured overhead differences. At the measured DES cost it projects up to
$\ProjectionMax\times$ speed-up, but does not execute DES-based search or establish
its exactness. Industrial evaluation costs can be substantially higher
\cite{saadatmand2025compdse}; transfer of the numerical speed-ups nevertheless requires
new measurements and bounds appropriate to that evaluator.

Conversely, the method offers little when the space is small, the constraints are loose,
the evaluation is nearly free, or --- most importantly --- when no admissible bound can be
constructed for the objectives of interest. If a bound cannot be proven admissible it must
not be used: the guarantee of Theorem~\ref{thm:exact} rests entirely on that property, and
an inadmissible bound silently converts an exact method into an approximate one without
any indication in the output.

\subsection{Relation to symbolic exact exploration}
Symbolic approaches~\cite{neubauer2018exact} prune partial designs with a constraint
solver and are, in principle, more general than closed-form bounds: any constraint the
solver can express can participate in the pruning. The assignment-monotone construction of the cited method does not directly
supply our queueing bounds. Our construction instead uses objective-specific
relaxations. These bounds preserve the analytical front, although the default order activates the exact queueing term only at full assignment. The order ablation tests earlier activation. Our finite-table Z3 baseline is exact but builds all evaluator values first; a compact solver front end with partial-design propagation remains an untested alternative.

\subsection{Relation to evaluator-centric methodologies}
CompDSE~\cite{saadatmand2025compdse} and HCA-DSE answer different questions. The former
asks how a design point of an industrial dCPS can be modelled and evaluated faithfully;
the latter asks which design points need not be evaluated at all. Their combination is
conditional: analytical bounds would need to be proved admissible for the
trace-derived simulator itself. Admissibility for the analytical model and good
empirical rank agreement do not establish that condition. The crossover analysis
is a cost projection and supplies no simulator-level exactness guarantee.

\subsection{Accuracy versus ordering}
The independent validation assesses both prediction error and ordering. Within-workload
agreement with measured mean-latency order is $\WithinAgreementRange\%$ on pairs
separated by more than $25\%$ in predicted latency; the pooled value is
$\PooledAgreement\%$. Pooling can benefit from workload separation and should not be
read as evidence that every local design choice is reliable. Repeated runs also expose
instability that a fixed analytical objective cannot represent.

The theorem guarantees the Pareto front of the analytical model. It does not imply that
this front equals the deployment or DES front. Analytical exploration can select
candidates for higher-fidelity reassessment, but discarding all other candidates preserves
that evaluator's front only if the pruning bounds are also proved admissible for it.
The measured-mean and measured-$p_{50}$ comparisons are reported separately because
an expected-latency approximation does not directly predict a latency quantile.
